% ---
% id: statistics/lem:freq-in-simplex
% kind: lem
% title: "Frequency vectors in the probability simplex"
% summary: >-
%   Every frequency vector has non-negative coordinates summing to one, since
%   each round contributes exactly one outcome. Intersecting a constraint set
%   with the probability simplex therefore never changes the probability that a
%   run passes the test, so simplex constraints cost no completeness.
% prominence: standard
% status: accepted
% uses: [def:freq]
% assumes: []
% refs: []
% provenance: newly-derived
% lean:
% ---

\begin{lemma}\label{lem:freq-in-simplex}
Let $n\ge1$, and write $\simplex{\Chatset}$ for the set of
$\mathbf{q}\in\mathbb{R}^{\Chatset}$ with non-negative coordinates summing
to one.
\begin{enumerate}[label=(\roman*),leftmargin=2.2em]
\item For every outcome string $\hat c_1^n\in\Chatset^{\,n}$, every count
$\Ncount{\hat c}(\hat c_1^n)$ of Definition~\ref{def:freq} is a non-negative
integer, at most $n$ of them are non-zero, and
$\sum_{\hat c\in\Chatset}\Ncount{\hat c}(\hat c_1^n)=n$. Hence
$\operatorname{freq}_{\hat c_1^n}\in\simplex{\Chatset}$.
\item For every set $S\subseteq\mathbb{R}^{\Chatset}$ and every random string
$\Chat_1^n$ with values in $\Chatset^{\,n}$, the sets
$\{\operatorname{freq}_{\Chat_1^n}\in S\cap\simplex{\Chatset}\}$ and
$\{\operatorname{freq}_{\Chat_1^n}\in S\}$ of sample points are equal. In
particular, whenever either is an event,
\begin{equation}\label{eq:simplex-free}
  \Pr\big[\operatorname{freq}_{\Chat_1^n}\in S\cap\simplex{\Chatset}\big]
  =\Pr\big[\operatorname{freq}_{\Chat_1^n}\in S\big].
\end{equation}
\end{enumerate}
\end{lemma}

\begin{proof}
\emph{(i).} Fix $\hat c_1^n$. By \eqref{eq:freq},
$\Ncount{\hat c}(\hat c_1^n)$ is the cardinality of the set
$J_{\hat c}\coloneqq\{j\in\{1,\dots,n\}:\hat c_j=\hat c\}$, a non-negative
integer. Each round $j\in\{1,\dots,n\}$ lies in $J_{\hat c}$ for exactly one
$\hat c$, namely $\hat c=\hat c_j$, so the sets $J_{\hat c}$ for
$\hat c\in\Chatset$ are pairwise disjoint with union $\{1,\dots,n\}$. Hence
at most $n$ of them are non-empty, and summing their cardinalities gives
$\sum_{\hat c\in\Chatset}\Ncount{\hat c}(\hat c_1^n)=n$. Dividing by $n\ge1$,
the coordinates $\operatorname{freq}_{\hat c_1^n}(\hat c)
=\tfrac1n\Ncount{\hat c}(\hat c_1^n)$ are non-negative and sum to one, which
is $\operatorname{freq}_{\hat c_1^n}\in\simplex{\Chatset}$.

\emph{(ii).} The inclusion
$\{\operatorname{freq}_{\Chat_1^n}\in S\cap\simplex{\Chatset}\}
\subseteq\{\operatorname{freq}_{\Chat_1^n}\in S\}$ holds because
$S\cap\simplex{\Chatset}\subseteq S$. Conversely, at every sample point the
value of $\Chat_1^n$ is a string in $\Chatset^{\,n}$, so by (i) the frequency
vector there lies in $\simplex{\Chatset}$; if it also lies in $S$, it lies in
$S\cap\simplex{\Chatset}$. The two sets are therefore equal, and equal events
have equal probability, which is \eqref{eq:simplex-free}.
\end{proof}
